% This is a modified version of the tufte-latex book example in which the title page and the contents page resemble Tufte's VDQI book, using Kevin Godby's code from this thread at https://groups.google.com/forum/#!topic/tufte-latex/ujdzrktC1BQ.
%
%%%%%%%%%%%%%%%%%%%%%%%%%%%%%%%%%%%%%%%%%%%%%%%%%%%%%%%%%%%%%%%%%%%%%%
% How to use Overleaf: 
%
% You edit the source code here on the left, and the preview on the
% right shows you the result within a few seconds.
%
% Bookmark this page and share the URL with your co-authors. They can
% edit at the same time!
%
% You can upload figures, bibliographies, custom classes and
% styles using the files menu.
%
% If you're new to LaTeX, the wikibook is a great place to start:
% http://en.wikibooks.org/wiki/LaTeX
%
%%%%%%%%%%%%%%%%%%%%%%%%%%%%%%%%%%%%%%%%%%%%%%%%%%%%%%%%%%%%%%%%%%%%%%
%% Unfortunately for the contents to contain
%% the "Parts" lines successfully, hyperref
%% needs to be disabled.
\documentclass[nohyper,nobib]{tufte-book}
\usepackage{nameref}
% \hypersetup{colorlinks}% uncomment this line if you prefer colored hyperlinks (e.g., for onscreen viewing)

% \usepackage{hyphenat}
\usepackage{url}
\usepackage[backend=biber, natbib=true, style=numeric]{biblatex}
\addbibresource{fsf-hyperspectral.bib}
\usepackage{xargs}
\renewcommandx{\cite}[3][1={0pt},2={}]{\sidenote[][#1]{\fullcite[#2]{#3}}}

%%
% Book metadata
\title{Calibration and \\ Characterisation of \\ Hyperspectral \\ Imagers}
\date{NERC FSF Protocols v 0.9}
%\author[The Tufte-LaTeX Developers (with modifications)]{The Tufte-LaTeX\ Developers (with modifications)}
%\publisher{NERC Field Spectroscopy Facility}

%%
% If they're installed, use Bergamo and Chantilly from www.fontsite.com.
% They're clones of Bembo and Gill Sans, respectively.
%\IfFileExists{bergamo.sty}{\usepackage[osf]{bergamo}}{}% Bembo
%\IfFileExists{chantill.sty}{\usepackage{chantill}}{}% Gill Sans

%\usepackage{microtype}

%%
% Just some sample text
\usepackage{lipsum}

%%
% For nicely typeset tabular material
\usepackage{booktabs}

\usepackage{xcolor}


%%
% For graphics / images
\usepackage{graphicx}
\setkeys{Gin}{width=\linewidth,totalheight=\textheight,keepaspectratio}
\graphicspath{{graphics/}}

% The fancyvrb package lets us customize the formatting of verbatim
% environments.  We use a slightly smaller font.
\usepackage{fancyvrb}
\fvset{fontsize=\normalsize}

%%
% Prints argument within hanging parentheses (i.e., parentheses that take
% up no horizontal space).  Useful in tabular environments.
\newcommand{\hangp}[1]{\makebox[0pt][r]{(}#1\makebox[0pt][l]{)}}

%%
% Prints an asterisk that takes up no horizontal space.
% Useful in tabular environments.
\newcommand{\hangstar}{\makebox[0pt][l]{*}}

%%
% Prints a trailing space in a smart way.
\usepackage{xspace}

%%
% Some shortcuts for Tufte's book titles.  The lowercase commands will
% produce the initials of the book title in italics.  The all-caps commands
% will print out the full title of the book in italics.
\newcommand{\vdqi}{\textit{VDQI}\xspace}
\newcommand{\ei}{\textit{EI}\xspace}
\newcommand{\ve}{\textit{VE}\xspace}
\newcommand{\be}{\textit{BE}\xspace}
\newcommand{\VDQI}{\textit{The Visual Display of Quantitative Information}\xspace}
\newcommand{\EI}{\textit{Envisioning Information}\xspace}
\newcommand{\VE}{\textit{Visual Explanations}\xspace}
\newcommand{\BE}{\textit{Beautiful Evidence}\xspace}

\newcommand{\TL}{Tufte-\LaTeX\xspace}

% Prints the month name (e.g., January) and the year (e.g., 2008)
\newcommand{\monthyear}{%
  \ifcase\month\or January\or February\or March\or April\or May\or June\or
  July\or August\or September\or October\or November\or
  December\fi\space\number\year
}


% Prints an epigraph and speaker in sans serif, all-caps type.
\newcommand{\openepigraph}[2]{%
  %\sffamily\fontsize{14}{16}\selectfont
  \begin{fullwidth}
  \sffamily\large
  \begin{doublespace}
  \noindent\allcaps{#1}\\% epigraph
  \noindent\allcaps{#2}% author
  \end{doublespace}
  \end{fullwidth}
}

% Inserts a blank page
\newcommand{\blankpage}{\newpage\hbox{}\thispagestyle{empty}\newpage}

\usepackage{units}
\usepackage{amsmath}

% Typesets the font size, leading, and measure in the form of 10/12x26 pc.
\newcommand{\measure}[3]{#1/#2$\times$\unit[#3]{pc}}

% Macros for typesetting the documentation
\newcommand{\hlred}[1]{\textcolor{Maroon}{#1}}% prints in red
\newcommand{\hangleft}[1]{\makebox[0pt][r]{#1}}
\newcommand{\hairsp}{\hspace{1pt}}% hair space
\newcommand{\hquad}{\hskip0.5em\relax}% half quad space
\newcommand{\TODO}{\textcolor{red}{\bf TODO!}\xspace}
\newcommand{\ie}{\textit{i.\hairsp{}e.}\xspace}
\newcommand{\eg}{\textit{e.\hairsp{}g.}\xspace}
\newcommand{\na}{\quad--}% used in tables for N/A cells
\providecommand{\XeLaTeX}{X\lower.5ex\hbox{\kern-0.15em\reflectbox{E}}\kern-0.1em\LaTeX}
\newcommand{\tXeLaTeX}{\XeLaTeX\index{XeLaTeX@\protect\XeLaTeX}}
% \index{\texttt{\textbackslash xyz}@\hangleft{\texttt{\textbackslash}}\texttt{xyz}}
\newcommand{\tuftebs}{\symbol{'134}}% a backslash in tt type in OT1/T1
\newcommand{\doccmdnoindex}[2][]{\texttt{\tuftebs#2}}% command name -- adds backslash automatically (and doesn't add cmd to the index)
\newcommand{\doccmddef}[2][]{%
  \hlred{\texttt{\tuftebs#2}}\label{cmd:#2}%
  \ifthenelse{\isempty{#1}}%
    {% add the command to the index
      \index{#2 command@\protect\hangleft{\texttt{\tuftebs}}\texttt{#2}}% command name
    }%
    {% add the command and package to the index
      \index{#2 command@\protect\hangleft{\texttt{\tuftebs}}\texttt{#2} (\texttt{#1} package)}% command name
      \index{#1 package@\texttt{#1} package}\index{packages!#1@\texttt{#1}}% package name
    }%
}% command name -- adds backslash automatically
\newcommand{\doccmd}[2][]{%
  \texttt{\tuftebs#2}%
  \ifthenelse{\isempty{#1}}%
    {% add the command to the index
      \index{#2 command@\protect\hangleft{\texttt{\tuftebs}}\texttt{#2}}% command name
    }%
    {% add the command and package to the index
      \index{#2 command@\protect\hangleft{\texttt{\tuftebs}}\texttt{#2} (\texttt{#1} package)}% command name
      \index{#1 package@\texttt{#1} package}\index{packages!#1@\texttt{#1}}% package name
    }%
}% command name -- adds backslash automatically
\newcommand{\docopt}[1]{\ensuremath{\langle}\textrm{\textit{#1}}\ensuremath{\rangle}}% optional command argument
\newcommand{\docarg}[1]{\textrm{\textit{#1}}}% (required) command argument
\newenvironment{docspec}{\begin{quotation}\ttfamily\parskip0pt\parindent0pt\ignorespaces}{\end{quotation}}% command specification environment
\newcommand{\docenv}[1]{\texttt{#1}\index{#1 environment@\texttt{#1} environment}\index{environments!#1@\texttt{#1}}}% environment name
\newcommand{\docenvdef}[1]{\hlred{\texttt{#1}}\label{env:#1}\index{#1 environment@\texttt{#1} environment}\index{environments!#1@\texttt{#1}}}% environment name
\newcommand{\docpkg}[1]{\texttt{#1}\index{#1 package@\texttt{#1} package}\index{packages!#1@\texttt{#1}}}% package name
\newcommand{\doccls}[1]{\texttt{#1}}% document class name
\newcommand{\docclsopt}[1]{\texttt{#1}\index{#1 class option@\texttt{#1} class option}\index{class options!#1@\texttt{#1}}}% document class option name
\newcommand{\docclsoptdef}[1]{\hlred{\texttt{#1}}\label{clsopt:#1}\index{#1 class option@\texttt{#1} class option}\index{class options!#1@\texttt{#1}}}% document class option name defined
\newcommand{\docmsg}[2]{\bigskip\begin{fullwidth}\noindent\ttfamily#1\end{fullwidth}\medskip\par\noindent#2}
\newcommand{\docfilehook}[2]{\texttt{#1}\index{file hooks!#2}\index{#1@\texttt{#1}}}
\newcommand{\doccounter}[1]{\texttt{#1}\index{#1 counter@\texttt{#1} counter}}

% Generates the index
\usepackage{makeidx}
\makeindex

%%%% Kevin Godny's code for title page and contents from https://groups.google.com/forum/#!topic/tufte-latex/ujdzrktC1BQ
\makeatletter
\renewcommand{\maketitlepage}{%
\begingroup%
\setlength{\parindent}{0pt}

{\fontsize{24}{24}\selectfont\textit{\@author}\par}

\vspace{1.75in}{\fontsize{36}{54}\selectfont\@title\par}

\vspace{0.5in}{\fontsize{14}{14}\selectfont\textsf{\smallcaps{\@date}}\par}

\vfill{\fontsize{14}{14}\selectfont\textit{\@publisher}\par}

\thispagestyle{empty}
\endgroup
}
\makeatother

\titlecontents{part}%
    [0pt]% distance from left margin
    {\addvspace{0.25\baselineskip}}% above (global formatting of entry)
    {\allcaps{Part~\thecontentslabel}\allcaps}% before w/ label (label = ``Part I'')
    {\allcaps{Part~\thecontentslabel}\allcaps}% before w/o label
    {}% filler and page (leaders and page num)
    [\vspace*{0.5\baselineskip}]% after

\titlecontents{chapter}%
    [4em]% distance from left margin
    {}% above (global formatting of entry)
    {\contentslabel{2em}\textit}% before w/ label (label = ``Chapter 1'')
    {\hspace{0em}\textit}% before w/o label
    {\qquad\thecontentspage}% filler and page (leaders and page num)
    [\vspace*{0.5\baselineskip}]% after
%%%% End additional code by Kevin Godby

\begin{document}

% Front matter
\frontmatter

% r.1 blank page
% \blankpage

% v.2 epigraphs
% \newpage\thispagestyle{empty}
% \openepigraph{%
% The public is more familiar with bad design than good design.
% It is, in effect, conditioned to prefer bad design, 
% because that is what it lives with. 
% The new becomes threatening, the old reassuring.
% }{Paul Rand%, {\itshape Design, Form, and Chaos}
% }
% \vfill
% \openepigraph{%
% A designer knows that he has achieved perfection 
% not when there is nothing left to add, 
% but when there is nothing left to take away.
% }{Antoine de Saint-Exup\'{e}ry}
% \vfill
% \openepigraph{%
% \ldots the designer of a new system must not only be the implementor and the first 
% large-scale user; the designer should also write the first user manual\ldots 
% If I had not participated fully in all these activities, 
% literally hundreds of improvements would never have been made, 
% because I would never have thought of them or perceived 
% why they were important.
% }{Donald E. Knuth}


% r.3 full title page
\maketitle
\begin{figure*}[h]
  \includegraphics[width=7.99cm, height=2cm]{FSF_Logo_Dark_Small}
\end{figure*}


% v.4 copyright page
%\newpage
%\begin{fullwidth}
%~\vfill
%\thispagestyle{empty}
%\setlength{\parindent}{0pt}
%\setlength{\parskip}{\baselineskip}
%Copyright \copyright\ \the\year\ \thanklessauthor

%\par\smallcaps{Published by \thanklesspublisher}

%\par\smallcaps{tufte-latex.googlecode.com}

%\par Licensed under the Apache License, Version 2.0 (the ``License''); you may not
%use this file except in compliance with the License. You may obtain a copy
%of the License at \url{http://www.apache.org/licenses/LICENSE-2.0}. Unless
%required by applicable law or agreed to in writing, software distributed
%under the License is distributed on an \smallcaps{``AS IS'' BASIS, WITHOUT
%WARRANTIES OR CONDITIONS OF ANY KIND}, either express or implied. See the
%License for the specific language governing permissions and limitations
%under the License.\index{license}

%\par\textit{First printing, \monthyear}
%\end{fullwidth}

% r.5 contents
\tableofcontents

\listoffigures

\listoftables

% r.7 dedication
%\cleardoublepage
%~\vfill
%\begin{doublespace}
%\noindent\fontsize{18}{22}\selectfont\itshape
%\nohyphenation
%Dedicated to those who appreciate \LaTeX{} 
%and the work of \mbox{Edward R.~Tufte} 
%and \mbox{Donald E.~Knuth}.
%\end{doublespace}
%\vfill
%\vfill


% r.9 introduction
\cleardoublepage
\chapter*{Overview}

\newthought{This document provides} an overall protocol to follow for the characterisation and calibration of hyperspectral imagers. 

\newthought{This protocol provides a series of tests} which can be used to determine the characteristics of a hyperspectral imager (hereafter, the Unit Under Test, or UUT), and to conduct a spectral and radiometric calibration of the UUT. 
Additionally, this document provides tests to determine imperfections in UUT image quality, such as co-registration errors and development of a bad pixel map,  concluding with steps for an uncertainty analysis of the radiometric calibration of the UUT. 
Throughout, an example application of the tests provided is conducted using -- as the example UUT -- a Headwall Co-Aligned Hyperspectral Imager HP.
The Python code used to process the raw data into the metrics presented here is provided in the companion code repository that accompanies this document.

\newthought{The development of this protocol} is based on the work conducted as part of the IEEE P4001 Working Group, the "Standard for Characterization and Calibration of Ultraviolet through Shortwave Infrared (250 nm to 2500 nm) Hyperspectral Imaging Devices" \cite{Durell20019} \cite{Ientilucci2022}



%%
% Start the main matter (normal chapters)
\mainmatter

\part{Spectral Domain}

\chapter{Spectral Characterisation}
\label{ch:SpectralCharacterisation}

\section{Introduction and Assumptions}
\newthought{The spectral characterisation} tests are aimed at determining the shapes of the spectral response features of select bands of the imager, with the assumption made that characterising a small sub-section of the spatial-spectral array will hold for all bands of the imager. 
By shapes, the tests aim to provide a point spread function of the response of each band to incoming signal. From this, assuming a Gaussian distribution in the response of the band, characteristics of the spectral response of the imager can be determined, namely, the spectral interval and  Full Width Half Maximum (FWHM), the latter of which provides the spectral resolution \sidenote{The relationship between the effective spectral resolution, $D_{\lambda eff}$, and the the FWHM, $D_{\lambda FWHM}$, can be mathematically determined from the spectral response, $R(\lambda)$: 
\begin{equation}
	D_{\lambda eff} = \frac{\int [R(\lambda )d \lambda]^{2}} {\int R^{2}(\lambda)d \lambda}
\end{equation} 
Which is equal to 1.504 $\cdot$ $D_{\lambda FWHM}$ }. 

\newthought{As mentioned}, it is assumed that the spectral response function will follow a Gaussian distribution. If this assumption does not hold, then the FWHM can not be used as a value for the spectral resolution of the instrument. Further tests utilising the data gathered using the method presented below can be used to determine whether this assumption holds. For the purposes of this section, it is assumed that the SRF will follow a Gaussian distribution. 

\newthought{It is also assumed} that characterising a subset of the spatial-spectral array is sufficient for the purposes of characterising the entirety of the array, provided that the subset samples pixels which are evenly distributed across the spatial-spectral array. In this way, the time required for testing and data processing is significantly reduced. This assumption holds if it is also assumed that the spectral response feature varies gradually across the spectral range and imager FOV.  A 10 $\times$ 5 sampling strategy i.e. 10 spectral band measurements at 5 spatial channels) of equally spaced spatial-spectral positions that span the spectral range and  FOV of the imager has been proven sufficient for the purposes of this test, although for full characterisation, the entirety of the array should be sampled. 

\newthought{Finally, it is assumed} that the monochromator used in the test conforms to the following specifications:
\begin{itemize}
	\item The monochromator itself must be spectrally characterised (following NERC FSF test protocols for monochromator calibration)
	\item The angular subtense of the image at the camera must exceed the sensor FOV to ensure full illumination of the spatial-spectral array. 
	\item The optical bandpass of the monochromator is less than 20\% of the quoted spectral sampling interval of the UUT, and allows for sub-band wavelength control e.g. the wavelength can be set to within 10\% of the spectral sampling interval or less.
\end{itemize}

\section{NERC FSF Method as applied to Headwall Co-Aligned VNIR-SWIR Imager}

\begin{enumerate}
	\item Attach the 2$''$ dia. integrating sphere to the output port of the OL 750-M-D monochromator. 
	\item Securely attach the UUT to the optical bench, and position the imager so that the VNIR fore optic optical axis aligns with the 2$''$  dia. integrating sphere port. The UUT will need to be vertically mounted, pointing at nadir, to view the interior of the 2$''$ dia. integrating sphere.
	\item Connect the lamp monochromator cables with the power supply. Switch on the power and supply, push "Output On/Off", and adjust the output of the power supply to 17.500 V
	\item Switch on the OL 750-C controller. 
	\item Choose an appropriate slit combination (refer to the OL 750-M-D monochromator additive sheet). 
	\item Start the OL 750 application software
	\item Select input and output slit widths which allow for an optical bandpass of about 20\% of the expected spectral sampling interval.
	\item Verify throughput of monochromator to integrating sphere.
	\item Start the UUT software.
	\item Wait for camera and monochromator to stabilise. 
	\item Select first monochromator centre wavelength, $\lambda_{N=1}$
	\item Verify that sufficient signal is present at the spectral band of interest. Vary the integration time of the sensor if required. Ensure that the full spatial channel is illuminated. 
	\item Reduce the monochromator wavelength so that the band of interest is no longer receiving signal. 
	\item Increase monochromator wavelength by ${\frac{1}{5}}^{th}$ to ${\frac{1}{10}}^{th}$ of one sampling interval and make note of the wavelength \sidenote{The decision on what fraction to set the sampling wavelength too will depend on the sampling interval and signal to noise ratio of the instrument. A smaller fraction of the sampling interval improves the determination of the shape of the SRF, but decreases the signal received by the UUT. A suggested rule of thumb is that the spectral source input should have a signal-to-noise ratio of 200 or better}.
	\item Cover entrance optic of UUT.  
	\item Name the datacube \textbf{NWW\_XX\_6090\_YY\_ZZ\_dark}. Here, WW is the test number; XX is the serial number of the UUT; YY is the monochromator centre wavelength; and ZZ is number of the interval conducted (e.g. if 20 points are to be taken across $\lambda_{N=1}$ , then 01 would be the first interval). 6090 is a numerical designator used by NERC FSF to signify that the datacube is a measurement used as part of a spectral characterisation test. 
	\item Take 100 dark frames.
	\item Remove cover.
	\item  Name the datacube \textbf{NWW\_XX\_6090\_YY\_ZZ\_light}. Take 100 light frames. 
	\item Repeat steps 14 to 19, incrementing ZZ, the interval number. Continue to sample across $\lambda_{N=1}$ until the output rises to a maximum value, decreases, and then no longer changes with monochromator wavelength. 
	\item Repeat steps 11 to 20, incrementing $\lambda_{N=1}$ e.g. $\lambda_{N=2}$, $\lambda_{N=3}$ ... until the desired number of spectral bands are sampled. 
	\item Repeat steps 1 to 21 for the SWIR array, altering step 2 to have the SWIR  fore optic optical axis aligns with the  2$''$ dia. integrating sphere port. 
\end{enumerate}

\section{Data Processing}

\begin{enumerate}
	\item Organize datacubes by monochromator centre wavelength, $\lambda_{N=n}$
	\item For each iteration, \textbf{NWW\_XX\_6090\_YY\_ZZ}, average the 100 dark and 100 light frames to arrive at a single averaged dark and light frame respectively. 
	\item Subtract the averaged dark frame from the averaged light frame to arrive at an averaged, dark subtracted light frame. Repeat for each \textbf{ZZ} iteration, arriving at \textbf{ZZ} number of averaged, dark subtracted light frames for each of the  $\lambda_{N=n}$ chosen wavelengths. 
	\item For each  $\lambda_{N=n}$, \textbf{ZZ} iteration, select an appropriate number of spatial channels to ensure that the changes in SRF response across the spatial dimension can be characterised. As an example, if using a 10 $\times$ 5 spectral band $\times$ spatial channel strategy, then 5 spatial channels should be chosen, and which are distributed evenly across the array. 
	\item For each of the selected spatial channels, plot the detector output as a function of monochromator wavelength after normalizing the peak (or area) to 1.0.
	\begin{figure}[h]
		\includegraphics{Ch2_S1_Characterisation.png}
		%  \checkparity This is an \pageparity\ page.%
 		 \caption[Example of SRF determination]{Example of SRF determination -- point spread function of, and Gaussian distribution fitting for, measurements taken at a select spectral band, for one spatial channel of the UUT. Plotted points were obtained by normalizing sequential responses from spatial channel no. 161 at spectral band no. 127 ( $\lambda_{N=4}$ in test case).}
 		 \label{fig:textfig}
 		 %\zsavepos{pos:textfig}
	\end{figure}
	\item Using this output, for each of the selected spatial channels for $\lambda_{N=n}$, determine if a Gaussian distribution curve can be fitted. If so, determine the FWHM of the fit, and from this, determine the effective spectral resolution. From this, determine the change in the FWHM across the spatial dimension for   $\lambda_{N=n}$.
	\item Repeat steps 2 to 6 for each  $\lambda_{N=n}$ e.g  $\lambda_{N=2}$,  $\lambda_{N=3}$ etc. From this, determine the change in the spectral dimension for the FWHM and effective spectral resolution. From this, an averaged FWHM and effective spectral resolution can be quoted for the entirety of the spatial-spectral array, provided a sufficient and evenly distributed sampling strategy has been performed. 
	\item Save data for analysis of the spectral co-registration error (SMILE) discussed in Part III, Errors of Co-registration and Imaging Imperfections. 
	\end{enumerate}


\chapter{Spectral Calibration}
\label{ch:SpectralCalibration}

\section{Introduction and Assumptions}
\newthought{The spectral calibration} aims to map each spectral band to a discrete wavelength number. This process is conducted by either determining empirically the centre wavelength for each spectral band, or more usually, determining the centre wavelength for select spectral bands, and developing a polynomial fit between those centre wavelengths which can then infer the centre wavelength of the remaining spectral bands. 

\newthought{While the tests} for spectral characterisation outlined above can be used for this procedure, it is often more expedient to use spectral line sources, such as cold cathode discharge lamps containing either a single gas (such as argon (Ar), xenon (Xe), krypton (Kr) etc), or a mixture of gases (such as mercury (Hg) and Ar). These tests are faster and require less setup, allowing for rapid spectral calibrations. The spectral characteristic tests are ideal for a first spectral calibration of an imager, but subsequent spectral calibrations can rely on the method outlined below. 

\newthought{It is assumed} that the spectral line sources used will, within their emission spectra, have emission lines which:
\begin{itemize}
	\item Provide throughput which exceeds the signal-to-noise threshold of the UUT
	\item Are discrete i.e. are not associated with other peaks in doublet configurations, and where the spectral distance between its centre wavelength and the centre wavelength of the succeeding and proceeding peaks exceeds, by a factor of three, the spectral resolution of the UUT \sidenote{The emission lines from commonly used spectral line sources, such as Hg, Ar, and Kr, are well known from both theory and measurements. An extensive database is maintained by NIST -https://www.nist.gov/pml/atomic-spectra-database}; and
	\item Have a narrow atomic line width, not exceeding  ${\frac{1}{10}}^{th}$ the spectral resolution of the UUT
\end{itemize}


\section{NERC FSF Method as applied to Headwall Co-Aligned VNIR-SWIR Imager}

\begin{enumerate}
	\item Securely attach the UUT to the optical bench,
	\item Position the 6$''$ dia. integrating sphere so that the output port aligns with the UUT's VNIR imager optical axis. 
	\item Start the UUT software.
	\item Insert the \textbf{Oriel 6035 Hg-Ar spectral line source} into the sphere's lamp port. 
	\item Set the power source to the appropriate current for the lamp, set to AC mode, and switch on.
	\item Wait a few moments, and switch to either DC+ or DC-. Allow 5 minutes to stabilize. If the lamp is to be on for a prolonged period of time, reverse DC polarity approximately every 5 minutes. 
	\item Verify that sufficient signal is present across the spectral range. Vary the integration time of the sensor if required. Ensure that the full spatial channel is illuminated. 
	\item Cover entrance optic of UUT.  
	\item Name the datacube \textbf{NYY\_ZZ\_6035\_dark}. Here, YY is the test number and ZZ is the serial number of the UUT. 6035 is a numerical designator used by NERC FSF to signify that the datacube is a measurement of an HgAr spectral line source as part of a spectral calibration test. 
	\item Take 100 dark frames
	\item Remove cover.
	\item Name the datacube \textbf{NYY\_ZZ\_6035\_light}. Take 1000 light frames.
	\item Switch off the line lamp. Wait 5 minutes to cool. 
	\item Remove the Oriel 6035 Hg-Ar spectral line source from the sphere lamp port. Insert the \textbf{Oriel 6030 Ar spectral line source}.  Repeat steps  5 to 12, renaming the datacube to reflect the new lamp (e.g.   \textbf{NYY\_ZZ\_6030\_dark},  \textbf{NYY\_ZZ\_6030\_light}.
	\item Remove the Oriel 6030 Ar spectral line source from the sphere lamp port. Insert the \textbf{Oriel 6031 Kr spectral line source}.  Repeat steps  5 to 12, renaming the datacube to reflect the new lamp (e.g.   \textbf{NYY\_ZZ\_6031\_dark},  \textbf{NYY\_ZZ\_6031\_light}.
	\item Position the 6$''$ dia. integrating sphere so that the output port aligns with the UUT's SWIR imager optical axis. Repeat steps 4 to 15.
	\item Remove the Oriel 6031 Kr spectral line source from the sphere lamp port. Insert the \textbf{Oriel 6033 Xe spectral line source}.  Repeat steps  5 to 12, renaming the datacube to reflect the new lamp (e.g.   \textbf{NYY\_ZZ\_6033\_dark},  \textbf{NYY\_ZZ\_6033\_light}.
	\item Switch off lamp and UUT. 
\end{enumerate}

\section{Data Processing}

\begin{enumerate}
	\item Organize datacubes by spectral line source number e.g. 6035, 6031 etc. 
	\item For each spectral line source, average the 100 dark and 1000 light frames to arrive at a single averaged dark and light frame respectively. 
	\item Subtract the averaged dark frame from the averaged light frame to arrive at an averaged, dark subtracted light frame. 
	\item Average the dark subtracted, average light frame across the spatial domain to give a single two column array of spectral band with measured signal.  \sidenote{Alternatively, take the central spatial channel to conduct the rest of the processing steps. A spectral calibration for each spatial channel can also be conducted, which is a methodology used for the determination of the spectral co-registration error, but as most imagers only accept a single set of spectral calibration coefficients, either the spatially averaged or single spatial channel method is used for the spectral calibration}. 
	\item For each of the frames processed in step 4, match the signal peaks to known emission lines for the respective line lamp (e.g. a signal peak detected at spectral band 17 is matched to the known emission line for HgAr at 404.66). If appropriate, conduct a Gaussian fit to the data to arrive at sub-spectral band number which is matched the emission line (e.g. after applying a Gaussian fit, the maximum signal peak for the HgAr emission line at 404.66 is found to be spectral band 16.7). 
	\item With the sub-spectral band peaks matched to emission lines for all spectral line sources, conduct a $3^{rd}$ order polynomial fit that relates band number ($\mathrm{N}$) to known wavelength ($\lambda(\mathrm{N}))$:
	\begin{equation}
		\lambda(\mathrm{N})=\mathrm{C_{0}}+\mathrm{C_{1}} \cdot(\mathrm{N})+\mathrm{C_{2}} \cdot(\mathrm{N})^2+\mathrm{C_{3}}\cdot(\mathrm{N})^3
	\end{equation}
	\item Save data for analysis of the spectral co-registration error (SMILE) discussed in Part III, Errors of Co-registration and Imaging Imperfections.  
\end{enumerate}
 

\part{Radiometric Domain}

\chapter{Introduction, Theory, and Assumptions}
\label{ch:radint}

\newthought{The radiometric characterisation and calibration} of hyperspectral imagers is, in theory, a simple process. For the calibration, the UUT is positioned so that it views a known radiance source, and measurements are taken of the UUT signal while viewing the source. After dark correction and normalisation, a spatial-spectral frame is generated of the measured signal. The expected radiance that should be measured by the UUT is determined by band-averaging the the known radiance using the measured spectral response features (generated via the process outlined in the first part of this document). Finally, this expected radiance is divided by the measured signal to arrive at a calibration coefficient for each spatial-spectral pixel of the detector:

\begin{equation}
	K_{(i, j)} = \frac{L_{(i,j)}}{S_{(i,j)}}
\end{equation}

where $K_{(i, j)}$ is the calibration coefficient for the array pixel situated at spatial channel $i$ and spectral band $j$, $L_{(i, j)}$ is the band averaged expected radiance for spatial-spectral pixel ($i$, $j$), and $S_{(i, j)}$ is the signal measured at spatial-spectral pixel ($i$, $j$). 

\newthought{Typically, the radiance source} used is an illuminated integrating sphere, with the UUT observing the interior. Alternatively, a reference plaque illuminated by an FEL standard can be used. In the methodology outlined below, the former setup is used. 

\newthought{The methodology outlined below} also allows for further characterisation tests to be conducted using the data gathered for the calibration.  Three characterisation metrics that can be determined by this approach, and which are related, are:
	\begin{itemize}
		\item The photon transfer curve
		\item A-star (A*)
		\item The noise equivalent spectral radiance (NESR)
	\end{itemize}
each of which is discussed in detail below.

\newthought{The development of a photon transfer curve} for the CCD of a UUT allows for the determination of the response of the CCD to uniform illumination. By plotting the temporal noise against the mean signal for each spatial-spectral pixel of the CCD array when viewing a uniform radiance source at various integration times of the UUT, various aspects of the CCD performance can be determined. These include the  \textbf{read noise}, which represents the random noise measured under totally dark conditions, the \textbf{photon shot noise}, which represents the response of the spatial-spectral pixel to the uniform source as integration time increases \cite{Janesick2007}. 

\newthought{By constructing the PTC} as a plot of the temporal variance against the mean dark-subtracted signal at different integration times, coefficients describing the non-linear behaviour of the CCD (\textbf{$\gamma$}) and the gain (\textbf{$G$}) in the in the photon shot noise region can be determined:

	\begin{equation}
		\sigma^{2} = \gamma \cdot x^{2} + G \cdot x + \epsilon
	\end{equation}

where $x$ is the signal and $\epsilon$ is the read noise as determined from the PTC. An example of such a plot is shown below:

	\begin{figure}[h]
		\includegraphics{3_PhotonTransferCurve.pdf}
		%  \checkparity This is an \pageparity\ page.%
 		 \caption[Example of PTC as variance versus mean signal]{Example of PTC as variance versus mean signal -- the temporal variance and mean dark-subtracted signal, averaged across spatial channels, for each spectral band of the UUT. Points are clustered by integration time.}
 		 \label{fig:textfig_1}
 		 %\zsavepos{pos:textfig_1}
	\end{figure}

\newthought{Using $\gamma$ and $G$}, the metric \textbf{A*} can be determined. A* is a parameter which is ideal for specifying the overall light collection of a camera \cite{Skauli2021}. A* is primarily a product of the geometrical factors representing light collection by the optics, and the factors representing the loss of light through the optics and detector. One means to express the A* for a UUT at a given wavelength is:

	\begin{equation}
		A* (\lambda) = \frac{k S_{DN}(\lambda)}{t_{int} \Delta\lambda L_{ph}(\lambda)}
	\end{equation}

where $k$ is inverse photon shot noise gain (i.e. the inverse of $G$), $S_{DN}(\lambda)$ is the raw signal at the specified wavelength, $t_{int}$ is the integration time, $\Delta\lambda$ is the bandwidth, and $L_{ph}(\lambda)$ is the photon spectral radiance at the specified wavelength, which can be calculated from the spectral radiance of the lamp in W multiplied by Planck's constant. 

\newthought{The noise equivalent spectral radiance} for a given wavelength can be determined if A* and the read noise are known:

	\begin{equation}
		NESR = \frac{hc}{\lambda} \frac{ \epsilon}{A* t_{int} \Delta\lambda}
	\end{equation}

where $h$ is Planck's constant, c is the speed of light, $lambda$ is the specified wavelength and $\epsilon$ is the read noise as calculated above. An example calculation of A* and NESR derived from A* is shown below:

	\begin{figure}[h]
		\includegraphics{5_Radiometric_Characterization.pdf}
		%  \checkparity This is an \pageparity\ page.%
 		 \caption[A* and NESR for an example UUT]{A* and NESR for an example UUT.}
 		 \label{fig:textfig_2}
 		 %\zsavepos{pos:textfig_1}
	\end{figure}

\newthought{It is assumed} that, if using the methodology below, the radiance sphere used will be well characterised (in terms of uniformity and geometrical properties) and calibrated (with regards to irradiance output of the lamp, the reflectance properties of the interior sphere wall, the at aperture radiance, and -- relatedly to the latter -- the output of the internal photodiode). 


\chapter{Methodology}
\label{ch:radiometricmethod}
\begin{enumerate}
	\item Bolt a 500 mm linear stage to the optical bench. Securely attach the UUT to the linear stage using an optical post attached to a post holder, to allow for both vertical and horizontal movement.
	\item Position the \textbf{Helios Dynamic Range 12" Radiance Source USLR-D12L-NMNN} radiance source on the optical table, and align the forward 12$''$ dia. output port with the UUT's VNIR imager optical axis. 
	\item Turn on the power supply for the sphere, and start the Helios software.
	\item Start the UUT software. 
	 \item Switch off overhead lights, and ensure sphere forward port is covered with the pinhole cover\sidenote{This cover has a pinhole situated at the exact centre of the port. This allows for the centring of the imager in both the horizontal and vertical spatial axes.} .
	\item  Select the HIS-075 calibration file, and zero the internal photometer. 
	\item In the Helios software, select the HIS-075 lamp, and set current and voltage to 6.13A and 12.00V. Allow the lamp to stabilise for 45 minutes. 
	\item Remove the optics cover of the UUT. Set the integration time in the UUT software to ensure maximum signal without saturation. 
	\item Gradually alter the height of the UUT using the optical post holder, until the central pixel of the central spatial channel is illuminated. Fix this position. 
	\item Remove the pinhole cover. 
	\item Gradually alter the horizontal distance between the UUT optics and the forward port of the sphere until the entire iFOV is illuminated. During this process, maintain a check on the internal photometer reading.
		\begin{enumerate}
			\item If the reading of the photometer increases as the UUT is brought closer to the forward port, then retro-reflection is occurring. Decrease the horizontal distance between the UUT and the forward port.
			\item For certain UUTs, it may not be possible to ensure that the full swath is illuminated without incurring retro-reflection. In this case, it will be required to proceed with the rest of the methodology by taking segmented measurements i.e. illuminating the first X number of spatial channels, then moving the sphere or UUT to illuminate the next X number of spatial channels, remembering to include an overlap region between segments. 
		\end{enumerate}
	\item Once the UUT is vertically and horizontally aligned without causing retro-reflection, fix position, and record distance from the sphere port. 
	\item Note the internal photometer reading. 
	\item Set the integration time of the UUT to where it is close to saturation. Note this integration time as $int_{max}$.
	\item Set the integration time of the UUT to as close to where minimum signal is achieved. Note this integration time as $int_{min}$
	\item Between $int_{min}$ and $int_{max}$, choose evenly distributed integration times, $int_{n}$. $int_{min}$ is set as $int_{n = 0}$, and $int_{max}$ is set as $int_{n = X}$, where X is the total number of integration time intervals decided upon.  
	\item Starting with $int_{n = 0}$:
		\begin{enumerate}
			\item Set integration time of the UUT to the integration time.
			\item Place cover over the forward port of the sphere, and a cover over the UUT optics.
			\item Name the datacube \textbf{NXX\_YY\_9001\_ZZ\_dark}. Here, XX is the test number; YY is the serial number of the UUT; and ZZ is the integration time interval, e.g. for  $int_{0}$ , ZZ = 1. 9001is a numerical designator used by NERC FSF to signify that the datacube is part of the radiometric calibration and characterisation test. 
			\item Take 100 dark frames
			\item Remove cover of forward port and UUT.
			\item Note the internal photometer reading. Name the datacube \textbf{NXX\_YY\_9001\_ZZ\_light}.Take 1000 light frames. 
		\end{enumerate}
	\item Repeat for each X interval between $int_{n = 0}$ $int_{n = X}$.
	\item Switch off the UUT and sphere. Allow the sphere 45 minutes to cool down.
	\item Repeat steps 1 to 19 for the SWIR imager of the UUT. 
\end{enumerate}

\chapter{Data Analysis}
\label{ch:rad-data}

\section{Photon Transfer Curve}
\begin{enumerate}
	\item Group \textbf{NXX\_YY\_9001} files by ZZ, the integration time of acquisition (e.g. 0.5 ms, 1.0 ms).
	\item For each integration time, average the 100 dark and 1000 light frames to arrive at a single averaged dark and light frame respectively.
	\item Subtract the averaged dark frame from the averaged light frame to arrive at an averaged, dark subtracted light frame, giving the mean signal $\bar{x}_{(i,j)}$ of each spatial-spectral pixel.
	\item Calculate the variance of each spatial-spectral pixel across the 1000 light frames, giving the temporal variance $\sigma^{2}_{(i,j)}$ of each spatial-spectral pixel.
	\item Average $\bar{x}_{(i,j)}$ and $\sigma^{2}_{(i,j)}$ across the spatial domain to give, for each spectral band, a single mean signal and a single variance.
	\item Combine the mean signal and variance for every spectral band into one array, grouped by integration time.
	\item Plot the variance against the mean signal, grouping by integration time.
	\item Conduct a $2^{nd}$ order polynomial fit of variance against mean signal to arrive at a fit following the below equation:
		\begin{equation}
			\sigma^{2} = \gamma \cdot x^{2} + G \cdot x + \epsilon
		\end{equation}
\end{enumerate}

\section{Radiometric Calibration}
	\begin{figure}[h]
		\includegraphics{3_L_Coeffs_Uncertainity.pdf}
		%  \checkparity This is an \pageparity\ page.%
 		 \caption[Example radiometric gains developed for UUT ]{Example radiometric gains developed for UUT. Each spatial-spectral pixel has its own associated radiometric gain. Likewise, uncertainties associated with these gains can be mapped similarly.}
 		 \label{fig:textfig_3}
 		 %\zsavepos{pos:textfig_1}
	\end{figure}
\begin{enumerate}
	\item Select a \textbf{NXX\_YY\_9001} group where the integration time provides a signal extending across the dynamic range of the UUT without saturation at any of the spectral bands (e.g. if 0.5 ms is the NXX\_YY\_9001 group with lowest measured signal, and 6.0 ms is NXX\_YY\_9001 group with the near saturation signal, then  the group at 4.0 ms, if using a 0.5 ms interval, would be an appropriate choice). 
	\item For the selected integration time group, average the 100 dark and 1000 light frames to arrive at a single averaged dark and light frame respectively.  Normalize each frame by the integration time i.e. divide the signal by the integration time to arrive a a counts per ms value e.g. if using the 4.0 ms frame, divide all values by 4 to arrive a counts per ms value. 
	\item Subtract the averaged, normalized dark frame from the averaged, normalized light frame to arrive at an averaged, dark subtracted, normalized light frame (hereafter, \textbf{NXX\_YY\_9001\_Radiance} .
	\item Import the known radiance of the sphere into a dataframe. 
	\item Using the spectral coefficients developed in the spectral calibration methodology, interpolate (band-average) the known radiance to arrive at an expected radiance value that should be measured across the spectral domain. 
	\item Divide this expected radiance by the measured signal across the spatial-spectral array to arrive at a 2-D map of radiometric coefficients for the UUT, as illustrated in Figure \ref{fig:textfig_3}
\end{enumerate}

\section{A* and NESR}
\begin{enumerate}
	\item Using the interpolated radiance value generated in part 5 of the Radiometric Calibration Data Analysis section, calculate the equivalent photon spectral radiance.
	\item Average across the spatial domain the \textbf{NXX\_YY\_9001\_Radiance} to arrive a two column array of averaged, dark corrected, normalized, and spatially averaged signal with associated spectral band. 
	\item Calculate A* using the following equation, where $k$ is the inverse of $G$ determined from the polynomial fit conducted in the Photon Transfer Curve, and $\epsilon$ is the read noise calculated, likewise, from the polynomial fit conducted in the Photon Transfer Curve :
		\begin{equation}
			A* (\lambda) = \frac{k S_{DN}(\lambda)}{t_{int} \Delta\lambda L_{ph}(\lambda)}
		\end{equation}
	\item Using the A* value, calculate NESR using the following equation:
	\begin{equation}
		NESR = \frac{hc}{\lambda} \frac{ \epsilon}{A* t_{int} \Delta\lambda}
	\end{equation}	
\end{enumerate}


\part{Errors of Co-registration and Imaging Imperfections}

\chapter{Errors of Co-registration}
\label{ch:coregerrors}

\newthought{When discussing} hyperspectral imagers, \textbf{co-registration} is defined as the alignment in position and shape of one response function with respect to another. This holds for both the spectral and spatial domain, i.e co-registration in the spectral domain is the alignment in the shape and centre wavelength of spectral bands across spatial pixels, while co-registration in the spatial domain is the alignment in the shape of the spatial point spread functions across spectral bands. 

	\begin{figure}[h]
		\includegraphics{Ch3_S1_CoRegErrors.png}
		%  \checkparity This is an \pageparity\ page.%
 		 \caption[Illustration of co-registration errors ]{Illustration of co-registration errors . Spectral smile leads to shifts in the reported centre wavelengths across the spatial array, while keystone leads to magnification across bands.}
 		 \label{fig:textfig_4}
 		 %\zsavepos{pos:textfig_4}
	\end{figure}

\newthought{An error in co-registration} is a deviation from the alignment in the position and shape of one response function with respect to another. When this occurs in the spectral domain, this is known as a \textbf{spectral smile}. It manifests as a centre wavelength shift, so that the wavelength of a spectral band measured at one end of the spatial array differs from the centre wavelength as measured at the centre or extreme of the array. For example, a pixel at spatial channel 1 may report a central wavelength of 400 nm for spectral band 40. The pixel at spatial channel 120, however, may report a central wavelength of 402 nm for the same spectral band, a 2 nm deviation in alignment. Spectral smile is an error associated with the line imaging spectrometer. 

\newthought{For the spatial domain}, an error in co-registration is known as spatial keystone. Here, the spatial point spread function varies across the spatial dimension, leading to pixels at one spectral band reporting a greater or lesser point spread function to other pixels, leading to a band to band magnification effect, which can result in over or under sampling. Spatial keystone is an error associated with the front optics of a spectrometer.


\section{Determination of Spectral Smile}
Both the methodology presented in Part One for spectral characterisation and for the spectral calibration can be used for the determination of the spectral smile for the UUT.  By taking measurements across multiple spectral bands, the centre wavelength for each pixel within that spatial band can be determined individually. This can then be compared to the averaged centre wavelength for the band, allowing for deviations to be quantified. The methodology outlined below is based on the data collected from the spectral calibration methodology. 

\begin{enumerate}
	\item Organize datacubes by spectral line source number e.g. 6035, 6031 etc. 
	\item For each spectral line source, average the 100 dark and 1000 light frames to arrive at a single averaged dark and light frame respectively. 
	\item Subtract the averaged dark frame from the averaged light frame to arrive at an averaged, dark subtracted light frame. 
	\item Deviating from the spectral calibration data analysis, do not average over the spatial array. 
	\item Match the signal peaks to known emission lines for the respective line lamp for every spectral channel.
	\item Conduct Gaussian fitting for every spatial channel to arrive at centre wavelengths for each spatial pixel per band, $\lambda{N}$
	\item For each peak identified, determine the offset of every value of $\lambda{c}$ from either the centre wavelength of the central spatial pixel ($\lambda{\#}$), or from the centre wavelength of the average of the spatial channels ($\lambda{\bar{x}}$)
	\item Conduct for all spectral line lamps. 
\end{enumerate}

An example for spectral smile determination for a UUT is shown below:

	\begin{figure}[h]
		\includegraphics{1_SMILE_Determination.pdf}
		%  \checkparity This is an \pageparity\ page.%
 		 \caption[Illustration of spectral smile ]{Illustration of spectral smile, using the deviation of  $\lambda{c}$  from $\lambda{\#}$ method.}
 		 \label{fig:textfig_5}
 		 %\zsavepos{pos:textfig_4}
	\end{figure}

\section{Determination of Spatial Keystone}
\newthought{A methodology for determining} spatial keystone is in development at NERC FSF. It utilizes a geometric control point methodology, where a grid pattern of pinholes at set spatial intervals is projected onto the imager array, creating a known interference pattern \cite{Bedard2012}. Future versions of this protocol document will present this methodology, including data analysis, further. 

\chapter{Imaging Imperfections}
\label{ch:imagingimp}
\newthought{Imaging imperfections} pertain to other metrics associated with the performance of the UUT which result in deviations from the ideal. These imperfections are associated with various aspects of the UUT design -- such as the CCD array or optics -- but can also be a result of instrument degradation. The prime imaging imperfection metrics to characterise include the linearity error, the effects of polarization, and \textbf{bad pixel mapping}. This protocol outlines a method for the latter only, but future versions of the protocol document will expand to include determining the former metrics. 

\section{Bad Pixel Mapping}
\newthought{A bad pixel} is defined as any pixel in the output image having atypical properties that do not conform to the specification of the UUT, i.e. which perform significantly different from the average pixel in the UUT array. The cause of these bad pixels can vary, but the impact is usually quantified by pixels which give no response when a response is expected (known as \textbf{cold pixels}),  or which give a noisy, inconsistent, or very high response (known as \textbf{hot pixels}). By taking the data generated during the radiometric calibration and characterisation methodology, it is possible to determine which pixels are differing from the mean response by a factor of four to five standard deviations. 

\begin{enumerate}
	\item  From the radiometric calibration methodology data analysis section, select a \textbf{NXX\_YY\_9001} group where the integration time provides a signal extending across the dynamic range of the UUT without saturation at any of the spectral bands (e.g. if 0.5 ms is the NXX\_YY\_9001 group with lowest measured signal, and 6.0 ms is NXX\_YY\_9001 group with the near saturation signal, then  the group at 4.0 ms, if using a 0.5 ms interval, would be an appropriate choice). 
	\item Calculate the standard deviation for each spatial-spectral pixel for the 1000 light frame datacube, generating a 2-D map of standard deviation per spatial-spectral pixel. 
	\item Calculate median standard deviation across the 2-D array. 
	\item Assuming a normal distribution in the standard deviation, calculate the thresholds that classify a cold pixel (which will have a standard deviation significantly lower than the median)  and a hot pixel (which will have a standard deviation significantly higher than the median). A good threshold is 4 to 5 standard deviations from the median standard deviation. 
	\item Apply thresholds to 2-D array, assigning a value of 1 to pixels that lie within the threshold range, and a value of 0 to those that do not.
	\item A 2-D map of pixels with either a value of 1 (normal) and 0 (bad) pixels will be generated, and can be applied post radiometric calibration to images to mask bad pixels. 
\end{enumerate}

	\begin{figure}[t]
		\includegraphics{BPM.pdf}
		%  \checkparity This is an \pageparity\ page.%
 		 \caption[Example of bad pixel map ]{Example of a bad pixel map. Pixels with a standard deviation (or in this case, uncertainty) significantly different from the median or mean pixel standard deviation (or uncertainty) are classified as bad, and given a value of 0. In this example, a few white pixels can be seen on the right hand-side image post BPM application -- these have been identified as bad pixels, and have been masked out.}
 		 \label{fig:textfig_6}
 		 %\zsavepos{pos:textfig_4}
	\end{figure}


\part{Uncertainty Analysis for Radiometric Calibration}


\chapter{Outline of an Uncertainty Estimation}
\label{ch:outlineuncertain}

\newthought{The methodology} for the estimation of uncertainty in the radiometric calibration of a UUT following the method provided in Part \ref{ch:radiometricmethod}  follows the procedure as outlined in the Guide to the Expression of Uncertainty in Measurement \cite{GUM2008}:

\begin{enumerate}
	\item \textbf{Determination of the Measurement Equation} - Express mathematically the relationship between the measurand $Y$ and the input quantities $X_{i}$ on which $Y$ depends: \( Y = f (X_{1}, X_{2}, ..., X_{N}) \). The function $f$ should contain every quantity, including all corrections and correction factors, that can contribute a significant component of uncertainty to the result of the
measurement 
	\item \textbf{Determination of Input Quantities} - Determine $x_{i}$, the estimated value of input quantity $X_{i}$, either on the basis of the statistical analysis of series of observations or by other means
	\item  \textbf{Determination of the Standard Uncertainty of Input Quantities} - Evaluate the standard uncertainty $u_{(x_{i})}$ of each input estimate $x_{i}$.
	\item \textbf{Evaluation of Covariance} - Evaluate the covariances associated with any input estimates that are correlated
	\item \textbf{Calculation of Measurement Result} - Calculate the result of the measurement, that is, the estimate $y$ of the measurand $Y$, from the functional
relationship $f$ using for the input quantities $X_{i}$ the estimates $x_{i}$ obtained in step 2.
	\item \textbf{Determination of combined uncertainties} - Determine the combined standard uncertainty $u_{c}$($y$) of the measurement result $y$ from the standard
uncertainties and covariances associated with the input estimate.
	\item \textbf{Expansion of uncertainties} - if necessary to provide an uncertainty range $U$, multiply the combined standard uncertainty $u_{c}$($y$) by a coverage factor $k$, typically in the range 2 to 3, to obtain \( U = k u_{c}(y) \)
	\item \textbf{Uncertainty Estimate} - Report the result of the measurement $y$ together with its combined standard uncertainty $u_{c}$($y$) or expanded uncertainty $U$
\end{enumerate}

\newthought{The following chapter} outlines the application of this methodology to determining the uncertainties associated with the radiometric calibration of a hyperspectral imager. 

\chapter{An Example Uncertainty Analysis for UUT}
\label{ch:uncertainityanalysis}

\section{Determination of the Measurement Equation}

\newthought{The \textbf{measurement equation}} is the expanded form of the \textbf{calculation equation}, which is the mathematical function relating the measurand to the input quantities. The measurement equation, as per GUM, must include all "corrections and correction factors that can contribute a significant component of uncertainty to the result of the measurement".  Stating first the calculation equation used for the radiometric calibration of the UUT, which is taken from  Part \ref{ch:radiometricmethod} :

\begin{equation}
	K_{(i, j)} = \frac{L_{(i,j)}}{S_{(i,j)}}
\end{equation}

where $K_{(i, j)}$, which is assigned as the measurand $Y$, is the radiometric calibration coefficient for a pixel at spatial channel $i$ at spectral band $j$. The input quantities $X_{i}$ for $Y$ are $L_{(i,j)}$,  the expected radiance based on the known, at aperture radiance of the sphere band averaged to the spectral response of the UUT; and $S_{(i,j)}$, the measured signal at spatial channel $i$ and spectral band $j$ of the UUT. For $S_{(i,j)}$, the calculation equation should be expanded to note that it is a product of the average collected light frames, $S_{light(i,j)}$ minus the average collected dark frames $S_{dark(i,j)}$, so that the calculation equation becomes:

\begin{equation}
	K_{(i, j)} = \frac{L_{(i,j)}}{S_{light(i,j)} - S_{dark(i,j)}}
\end{equation}

\newthought{Hidden assumption terms} associated with the input quantities must also be included to fully expanded the calculation equation into a measurement equation. For this uncertainty analysis, the hidden assumption terms (given the symbol $K$) include the drift in the irradiance output of the HIS-075 lamp within the radiance sphere $K_{drift}$, and the sphere uniformity $K_{uniformity}$. With these assumption terms considered, the final measurement equation is given as:

\begin{equation}
	K_{(i, j)} = \frac{L_{(i,j)}K_{drift}K_{uniformity}}{S_{light(i,j)} - S_{dark(i,j)}}
	\label{eq:measurement}
\end{equation}

\section{Determination of Input Quantities}

In total, five input quantities must be measured:
\begin{itemize}
	\item $L_{(i,j)}$ -- this is derived by band-averaging the known radiance of the sphere. The known radiance of the sphere is determined from the sphere calibration, which includes calibration of the HIS-075 lamp and the characterisation of the reflectance and geometrical properties of the interior sphere wall. The sphere calibration in itself requires an uncertainty estimation. This is provided by the sphere manufacturer. Additionally, in the process of band averaging the known radiance to an expected radiance, the spectral coefficients for the UUT are required, and consequently, the uncertainty associated with those coefficients will be necessary to calculate the uncertainty associated with  $L_{(i,j)}$.
	\item $K_{drift}$ -- this can be determined from the sphere internal photometer reading and the lifetime of the lamp. 
	\item $K_{uniformity}$ -- this describes the variation in the uniformity of the sphere, and is characterised during sphere calibration. In this case, the sphere uniformity is provided by the sphere manufacturer. 
	\item $S_{light(i,j)}$ -- this is the measured light signal of the UUT when observing the illuminated sphere interior. 
	\item $S_{dark(i,j)}$ -- this is the measured signal of the UUT when the optic is covered and all sources of incoming light are removed. 
\end{itemize}

\section{Determination of the Standard Uncertainty of Input Quantities}

\newthought{The following sections} will present the data from a VNIR imager calibration to illustrate the steps taken. Copies of the data are provided with this protocol, but for brevity, a single spatial channel, \# 320, spanning 35 spectral bands is presented. 

\newthought{The standard deviation} for both $S_{light(i,j)}$ and $S_{dark(i,j)}$ is determined from the data collected in the methodology present in  Part \ref{ch:radiometricmethod}. In Table \ref{tab:my-table-1} is the standard deviation, $\sigma$, for  $S_{light(320, 0-35)}$ and $S_{dark(320, 0-35)}$. Because $S_{light(i,j)}$ and $S_{dark(i,j)}$ are the means of $N_{light} = 1000$ and $N_{dark} = 100$ frames respectively, their standard uncertainties are the standard deviations of those means:

\begin{equation}
	u\left(S_{light(i,j)}\right) = \frac{\sigma_{light(i,j)}}{\sqrt{N_{light}}}, \qquad u\left(S_{dark(i,j)}\right) = \frac{\sigma_{dark(i,j)}}{\sqrt{N_{dark}}}
\end{equation}

% Please add the following required packages to your document preamble:
% \usepackage{booktabs}
\begin{table}[h]
\centering
\begin{tabular}{@{}lll@{}}
\toprule
Wavelength & $\sigma$ $S_{dark(320, 0-35)}$ & $\sigma$ $S_{light(320, 0-35)}$ \\ \midrule
400.96     & 9.73E-01 & 1.30E+00  \\
403.21     & 9.63E-01 & 1.42E+00  \\
405.47     & 9.16E-01 & 1.43E+00  \\
407.72     & 9.57E-01 & 1.50E+00  \\
409.98     & 9.72E-01 & 1.55E+00  \\
412.23     & 9.40E-01 & 1.74E+00  \\
414.49     & 9.50E-01 & 1.64E+00  \\
416.74     & 9.55E-01 & 1.77E+00  \\
419        & 9.03E-01 & 1.75E+00  \\
421.25     & 9.38E-01 & 1.83E+00  \\
423.51     & 9.46E-01 & 1.92E+00  \\
425.76     & 9.94E-01 & 1.88E+00  \\
428.02     & 9.18E-01 & 2.01E+00  \\
430.27     & 9.42E-01 & 2.06E+00  \\
432.53     & 9.38E-01 & 2.14E+00  \\
434.79     & 9.06E-01 & 2.10E+00  \\
437.04     & 9.39E-01 & 2.18E+00  \\
439.3      & 9.29E-01 & 2.23E+00  \\
441.55     & 9.13E-01 & 2.33E+00  \\
443.81     & 1.04E+00 & 2.27E+00  \\
446.07     & 9.83E-01 & 2.34E+00  \\
448.32     & 9.76E-01 & 2.49E+00  \\
450.58     & 9.42E-01 & 2.69E+00  \\
452.83     & 9.57E-01 & 2.51E+00  \\
455.09     & 9.68E-01 & 2.62E+00  \\
457.35     & 9.62E-01 & 2.59E+00  \\
459.6      & 9.50E-01 & 2.78E+00  \\
461.86     & 9.58E-01 & 2.77E+00  \\
464.12     & 9.51E-01 & 2.98E+00  \\
466.37     & 9.82E-01 & 2.90E+00  \\
468.63     & 9.28E-01 & 3.02E+00  \\
470.89     & 9.53E-01 & 2.94E+00  \\
473.14     & 9.82E-01 & 3.14E+00  \\
475.4      & 9.70E-01 & 3.18E+00  \\
477.66     & 9.31E-01 & 3.30E+00  \\ \bottomrule
\end{tabular}
\caption{Standard deviation in measured dark and light frames across spectral range for spatial channel \# 320 for example UUT.}
\label{tab:my-table-1}
\end{table}

As the uncertainty associated with $L_{(i,j)}$ is quantified during sphere calibration and subsequently expressed as an uncertainty, no associated standard deviation is associated with $L_{(i,j)}$. Where the sphere calibration certificate quotes an expanded uncertainty, it is divided by the stated coverage factor (typically $k = 2$) to give the standard uncertainty $u(L_{(i,j)})$. For completeness sake, the interpolated values for $L_{(320,0-35)}$ and interpolated values for $u_{L}$ are included in Table \ref{tab:my-table-2}. The determination of $K_{drift}$ is dependent on the internal photometer reading and lamp lifetime. From experiments, a flat 0.2\% uncertainty is estimated. For $K_{uniformity}$ , deviations in internal sphere uniformity is measured using an imaging photometer viewing the sphere interior, and calculating the standard deviation from these measurements. From this, a flat 0.4\% uncertainty is estimated. Both $K_{drift}$ and $K_{uniformity}$ take a nominal value of 1, so that they contribute to the uncertainty of $K_{(i,j)}$ without altering its value. 

% Please add the following required packages to your document preamble:
% \usepackage{booktabs}
% \usepackage[normalem]{ulem}
% \useunder{\uline}{\ul}{}
\begin{table}[]
\centering
\begin{tabular}{@{}lll@{}}
\toprule
Wavelength (nm) & $L_{(320,0-35)}$ (mW m$^{-2}$ sr$^{-1}$ nm$^{-1}$) & $u_{L}$ (mW m$^{-2}$ sr$^{-1}$ nm$^{-1}$) \\ \midrule
400.96          & 48.31             & 0.38              \\
403.21          & 51.3              & 0.38              \\
405.47          & 54.35             & 0.405             \\
407.72          & 58.17             & 0.425             \\
409.98          & 62.24             & 0.445             \\
412.23          & 66.17             & 0.47              \\
414.49          & 69.94             & 0.49              \\
416.74          & 73.81             & 0.51              \\
419             & 77.6              & 0.53              \\
421.25          & 81.54             & 0.56              \\
423.51          & 85.39             & 0.58              \\
425.76          & 89.3              & 0.605             \\
428.02          & 93.03             & 0.63              \\
430.27          & 96.75             & 0.65              \\
432.53          & 100.08            & 0.67              \\
434.79          & 103.3             & 0.68              \\
437.04          & 107.02            & 0.795             \\
439.3           & 110.88            & 0.815             \\
441.55          & 114.53            & 0.835             \\
443.81          & 118.29            & 0.85              \\
446.07          & 122.1             & 0.88              \\
448.32          & 126.02            & 0.895             \\
450.58          & 129.99            & 0.92              \\
452.83          & 133.95            & 0.94              \\
455.09          & 137.95            & 0.96              \\
457.35          & 141.83            & 0.98              \\
459.6           & 145.66            & 0.995             \\
461.86          & 149.46            & 1.02              \\
464.12          & 153.43            & 1.055             \\
466.37          & 157.77            & 1.08              \\ \bottomrule
\end{tabular}
\caption{Interpolated sphere radiance and interpolated uncertainty associated with sphere radiance for example UUT.}
\label{tab:my-table-2}
\end{table}


\section{Evaluation of Covariance}
\newthought{The combined standard uncertainty}, $u_{c}(y)$, is developed using a first-order Taylor series approximation, known as the law of propagation of uncertainty. Here, the standard uncertainties associated with the input parameters $x_{i}$ are propagated through:

\begin{equation}
u_c^2(y)=  \sum_{i=1}^N\left(\frac{\partial f}{\partial x_i}\right)^2 u^2\left(x_i\right) +2 \sum_{i=1}^{N-1} \sum_{j=i+1}^N \frac{\partial f}{\partial x_i} \frac{\partial f}{\partial x_j} u\left(x_i, x_j\right)
\end{equation}

where $ u\left(x_i, x_j\right)$  expresses the covariance between input parameters. If input parameters are independent and not correlated, as in this example, then the covariance term vanishes. Therefore, an analysis of the covariance between inputs is not required for this example. 

\section{Calculation of Measurement Result}

\newthought{With $L_{(i,j)}$, $S_{light(i,j)}$, and $S_{dark(i,j)}$} known, $K_{(i, j)}$ can be calculated for each spatial-spectral pixel. As an illustration, the values for $K_{(320, 0 -30)}$ are expressed in Table \ref{tab:my-table-3}.

% Please add the following required packages to your document preamble:
% \usepackage{booktabs}
\begin{table}[]
\centering
\begin{tabular}{@{}ll@{}}
\toprule
Wavelength (nm) & K        \\ \midrule
400.96          & 4.15E+00 \\
403.21          & 3.68E+00 \\
405.47          & 3.61E+00 \\
407.72          & 3.19E+00 \\
409.98          & 3.18E+00 \\
412.23          & 2.86E+00 \\
414.49          & 2.94E+00 \\
416.74          & 2.64E+00 \\
419             & 2.69E+00 \\
421.25          & 2.43E+00 \\
423.51          & 2.48E+00 \\
425.76          & 2.33E+00 \\
428.02          & 2.34E+00 \\
430.27          & 2.17E+00 \\
432.53          & 2.14E+00 \\
434.79          & 2.02E+00 \\
437.04          & 2.02E+00 \\
439.3           & 1.91E+00 \\
441.55          & 1.86E+00 \\
443.81          & 1.74E+00 \\
446.07          & 1.73E+00 \\
448.32          & 1.65E+00 \\
450.58          & 1.64E+00 \\
452.83          & 1.52E+00 \\
455.09          & 1.50E+00 \\
457.35          & 1.41E+00 \\
459.6           & 1.43E+00 \\
461.86          & 1.33E+00 \\
464.12          & 1.34E+00 \\
466.37          & 1.28E+00 \\ \bottomrule
\end{tabular}
\caption{Radiometric gain per spectral band for spatial channel \# 320 for example UUT.}
\label{tab:my-table-3}
\end{table}


\section{Determination of combined uncertainties}

\newthought{With the standard uncertainties} of $L_{(i,j)}$, $K_{drift}$, $K_{uniformity}$, $S_{light(i,j)}$, and $S_{dark(i,j)}$ known, they can be combined to provide a combined standard uncertainty for $K_{(i,j)}$:

\begin{equation}
\label{eq:combined}
\begin{aligned}
u_{c}^{2}\left(K_{(i,j)}\right) = {} & \left( \frac{\partial K_{(i,j)}}{\partial L_{(i,j)}} \right)^{2} u^{2}\left(L_{(i,j)}\right) \\
& + \left( \frac{\partial K_{(i,j)}}{\partial K_{drift}} \right)^{2} u^{2}\left(K_{drift}\right) \\
& + \left( \frac{\partial K_{(i,j)}}{\partial K_{uniformity}} \right)^{2} u^{2}\left(K_{uniformity}\right) \\
& + \left( \frac{\partial K_{(i,j)}}{\partial S_{light(i,j)}} \right)^{2} u^{2}\left(S_{light(i,j)}\right) \\
& + \left( \frac{\partial K_{(i,j)}}{\partial S_{dark(i,j)}} \right)^{2} u^{2}\left(S_{dark(i,j)}\right)
\end{aligned}
\end{equation}

Here, the partial derivative $\frac{\partial K_{(i,j)}}{\partial x}$ is the sensitivity coefficient associated with the standard uncertainty of input quantity $x$. The sensitivity coefficients can be determined in a number of ways. In this example, they were determined analytically by differentiating Eq. \ref{eq:measurement}, to provide the following:

\begin{table}[h]
\centering
\begin{tabular}{@{}ll@{}}
\toprule
Input Quantity $x$ & Sensitivity Coefficient $\partial K_{(i,j)} / \partial x$ \\ \midrule
 $L_{(i,j)}$ & $K_{(i,j)} / L_{(i,j)}$ \\
 $K_{drift}$ & $K_{(i,j)} / K_{drift}$ \\
 $K_{uniformity}$ & $K_{(i,j)} / K_{uniformity}$ \\
 $S_{light(i,j)}$ & $-K_{(i,j)} / \left(S_{light(i,j)} - S_{dark(i,j)}\right)$ \\
 $S_{dark(i,j)}$ & $+K_{(i,j)} / \left(S_{light(i,j)} - S_{dark(i,j)}\right)$ \\ \bottomrule
\end{tabular}
\caption{Sensitivity coefficients for example uncertainty analysis}
\label{tab:my-table-4}
\end{table}

Substituting these coefficients into Eq. \ref{eq:combined} and dividing by $K_{(i,j)}^{2}$ gives the combined uncertainty in relative terms, which is the form used in practice, as the sphere, drift, and uniformity uncertainties are quoted as percentages. Dropping the pixel indices for brevity:

\begin{equation}
\label{eq:relative}
\begin{aligned}
\frac{u_{c}^{2}(K)}{K^{2}} = {} & \frac{u^{2}(L)}{L^{2}} + \frac{u^{2}(K_{drift})}{K_{drift}^{2}} + \frac{u^{2}(K_{uniformity})}{K_{uniformity}^{2}} \\
& + \frac{u^{2}(S_{light}) + u^{2}(S_{dark})}{\left(S_{light} - S_{dark}\right)^{2}}
\end{aligned}
\end{equation}

Using Eq. \ref{eq:relative} with the standard uncertainties determined above, the combined standard uncertainty for each spatial-spectral pixel can be determined.

\section{Expansion of uncertainties}
A coverage factor of two was applied to the combined uncertainty to give an uncertainty range, $U_{K}$, for each spatial-spectral pixel. 

\section{Uncertainty Estimate}
A final uncertainty estimate, given as a relative uncertainty, was then determined for each spatial-spectral pixel. 

%%
% The back matter contains appendices, bibliographies, indices, glossaries, etc.

\backmatter

%\bibliography{fsf-hyperspectral}
%\bibliographystyle{plainnat}


\printindex

\end{document}
